\documentclass[10pt,a4paper]{amsart}
\usepackage[margin=19mm]{geometry}
\usepackage{amsmath,amssymb,mathtools}
\usepackage[scr=rsfs,cal=euler]{mathalfa}
\usepackage{xfrac}
\usepackage[unicode,hidelinks,bookmarks=false]{hyperref}
\newcommand{\ie}{i.e.\ }
\newcommand{\cf}{cf.\ }
\newcommand{\resp}{respectively\ }
\newcommand{\Subsection}{\subsection}
\providecommand{\nobreakdash}{\nobreak\hbox{-}}
\setlength{\emergencystretch}{2em}
\multlinegap0pt
\usepackage{amssymb}
\newtheorem{lemma}{Lemma}
\DeclareMathOperator{\Ann}{Ann}
\DeclareMathOperator{\CM}{CM}
\DeclareMathOperator{\End}{End}
\DeclareMathOperator{\Id}{Id}
\DeclareMathOperator{\Ker}{Ker}
\title{Orthogonal completion of algebraic systems}
\author{A. Yu. Golubkov}
\date{Source: arXiv:2609.01854v1 (CC BY 4.0)}
\begin{document}
\maketitle

\noindent\emph{Source and licence.} Orthogonal completion of algebraic systems. Version arXiv:2609.01854v1 (CC BY 4.0), distributed by arXiv under the Creative Commons Attribution 4.0 International licence, http://creativecommons.org/licenses/by/4.0/, as recorded in the arXiv licence metadata for this exact version and shown on the abstract page https://arxiv.org/abs/2609.01854v1. Full licence notice: \texttt{licenses/arxiv-2609.01854-CC-BY-4.0.txt}.\par\smallskip

\noindent\textit{Editorial scope.} Counted unit: the article's lemma lemma (source0.tex lines 1036--1044). The article's complete original proof of it is delivered with it (source0.tex lines 1046--1159, one \texttt{proof} environment of 114 lines). No mathematics is written, repaired or re-derived: the statement and the proof are the article's own lines, in the article's own notation, and no retained line is substituted or re-flowed.  No further source-local context is needed: every cross-reference of every retained line resolves inside the two delivered spans.  Nothing was omitted from any retained span, so the package carries no omission record.  No retained line contains a citation, so the wrapper carries no bibliography and no bibliography entry.  No retained line contains a figure environment, an image-inclusion command or drawing code, so no image is delivered and the package declares no asset dependency.  Editorial formatting only, in addition: an AMS \texttt{amsart} preamble that restates the article's own declarations and macros used by the retained lines, namely the environments \texttt{lemma} on separate counters, together with the standard packages those lines need.  The retained environments print in this document as Lemma 1 in order of appearance, whereas the article prints the same items with the numbering of its own full document.  The complete native source of the article as distributed in the arXiv e-print for this exact version is bundled unchanged as \texttt{sources/209147/source0.tex}.
\begin{lemma}
If $B$ is an ultrafilter in $B(R)$, ${\nu(B)\in B}$ or (and) there exist 
${k, n\geq 1}$ such that the elements of $M(O(P(R)))$ are the sums of $k$ 
products of at most $n$ operators $t_x$, 
${t=l, r}$, ${x\in O(P(R))}$, then $O(P(R))_B$ is an essential 
$M(O(P(R)))'_B$--submodule of $I(P(R))_B$ and 
${Z(M(O(P(R))_B)')\cong Z(M(O(P(R)))')_B\cong \CM(R)/P \CM(R)}$, 
${P=B(R)\setminus B}$.
\end{lemma}

\begin{proof}
The canonical epimorphism ${\psi_B: O(P(R))\longrightarrow O(P(R))_B}$ induces 
an epimorphism ${\psi_{M, B}: M(O(P(R)))'\longrightarrow M(O(P(R))_B)', 
\psi_{M, B} \Id_{O(P(R))}=\Id_{O(P(R))_B}, 
\psi_{M, B} t_x=t_{\psi_B x}}$ for all ${t=l, r}$, ${x\in O(P(R))}$, 
\begin{multline*}
\Ker \psi_{M, B}\ =\ \Ann_{M(O(P(R)))'} O(P(R))_B\ =
\\ 
\{\phi\in M(O(P(R)))'\mid (O(P(R)))\phi=O((O(P(R)))\phi)
\subseteq P O(P(R))\}\ =\ P M(O(P(R)))'
\end{multline*}
(there is ${\alpha\in B}$, ${\alpha (O(P(R)))\phi=\{0\}}$, ${\alpha \phi=0}$, 
${\phi\in P M(O(P(R)))'}$ (Proposition 2.4 and its proof)),
${M(O(P(R))_B)'\cong M(O(P(R)))'_B=M(O(P(R)))'/P M(O(P(R)))'}$. 
The structure of the $M(O(P(R)))'$--module is introduced on $I(P(R))$ in the 
proof of Proposition 2.4. When the second condition is satisfied for any dense 
orthogonal ${\{\xi_a\}_{a\in I}\subseteq B(R)}$, 
${\{x_a\}_{a\in I}\subseteq I(P(R))}$, 
${\{\phi_a\}_{a\in I}\subseteq M(O(P(R)))'}$, writing 
\[
\phi_a\ =\ \tau_a \Id_{O(P(R))}+
\sum_{j=1}^n \sum_{i=1}^k \sum_{T=(t^{(1)}, \ldots, t^{(j)})\in 
\{l, r\}^j} t^{(1)}_{x(a, i, T, 1)}\cdots t^{(j)}_{x(a, i, T, j)}
\]
for ${\tau_a\in \CM(R)}$, ${x(a, i, T, j)\in O(P(R)), a\in I}$, where the sum 
of ${1+k (2^{n+1}-2)}$ terms has at most ${k+1}$ non-zero terms, we get 
\[
{\sum_{a\in I}}^{\perp} \xi_a (x_a \phi_a)\ =\ 
\tau x+
\sum_{j=1}^n \sum_{i=1}^k \sum_{T=(t^{(1)}, \ldots, t^{(j)})\in 
\{l, r\}^j} x t^{(1)}_{x(i, T, 1)}\cdots t^{(j)}_{x(i, T, j)}\,,
\]
${\tau\in \CM(R)}$, ${\xi_a \tau=\xi_a \tau_a, a\in I}$ (Lemma 2.3), 
${x={\sum\limits_{a\in I}}^{\perp} \xi_a x_a}$, 
${x(i, T, l)={\sum\limits_{a\in I}}^{\perp} \xi_a x(a, i, T, l)}.$ 
In particular, it follows that ${x M(O(P(R)))'=O(x M(O(P(R)))')}$, 
\[
(S)[\phi, M(O(P(R)))]\ =\ 
\bigcup\limits_{s\in S} s [\phi, M(O(P(R)))]\ =\ O((S)[\phi, M(O(P(R)))])
\]
for all ${x\in I(P(R))}$, ${S=O(S)\subseteq I(P(R))}$, ${\phi\in M(O(P(R)))'}$.

If ${x\in I(P(R))}$, ${x M(O(P(R)))'\cap O(P(R))\subseteq P I(P(R))}$, then for 
${\beta=\nu(B)\in B}$ for the first condition and some ${\beta\in B}$ 
for the second ${\beta x M(O(P(R)))'\cap O(P(R))=\{0\}}$, since for the latter 
${x M(O(P(R)))'\cap O(P(R))=O(x M(O(P(R)))'\cap O(P(R)))}$, 
\[
\beta x M(O(P(R)))'\cap O(P(R))\subseteq \beta (x M(O(P(R)))'\cap O(P(R)))\ =\ 
\{0\}\,,
\]
(Proposition 2.4, p. 2, 6). In any case, ${\beta x=0}$, ${x\in P I(P(R))}$. 
Hence 
\[
(x M(O(P(R)))')_B\cap O(P(R))_B\ =\ (x+P I(P(R))) M(O(P(R)))'_B\cap O(P(R))_B
\ \ne\ \{0\}
\] 
for all ${x\notin P I(P(R))}$. If ${\phi\in M(O(P(R)))'}$ and 
${\psi_{M, B} \phi\in Z(M(O(P(R))_B)')}$, then  
\[
(O(P(R)))[\phi, M(O(P(R)))]\subseteq P O(P(R))\,,\quad 
\gamma (O(P(R)))[\phi, M(O(P(R)))]\ =\ 
\{0\}
\]
with ${\gamma=\nu(B)}$ for the first condition and some ${\gamma\in B}$ for 
the second, ${\gamma \phi\in Z(M(O(P(R)))')}$, 
\begin{multline*}
\phi+P M(O(P(R)))'\in (Z(M(O(P(R)))')+P M(O(P(R)))')/P M(O(P(R)))'\cong
\\
Z(M(O(P(R))))'_B\cong \psi_{M, B}(Z(M(O(P(R)))'))
\end{multline*}
(Proposition 2.4, p. 2, 6). Since 
${\psi_{M, B}(Z(M(O(P(R)))'))\subseteq Z(M(O(P(R))_B)')}$, 
\begin{multline*}
Z(M(O(P(R))_B)')\ =\ \psi_{M, B}(Z(M(O(P(R)))'))\ =
\\ 
\psi_{M, B}(\CM(R) \Id_{O(P(R))})\cong \CM(R)/P \CM(R)
\end{multline*}
(Proposition 2.4, p. 4).

For ${\nu(B)\in B}$, the essentiality of the $M(R)'_B$--submodule 
${R_B\subseteq P(R)_B}$ and $M(P(R))'_B$--submodule 
${P(R)_B\subseteq I(P(R))_B}$ together with   
\[
Z(M(R_B)') \cong Z(M(R)')_B\,,\quad 
Z(M(P(R)_B)')\cong Z(M(P(R))')_B\cong \CM(R)/P \CM(R)
\]
(${M(R)'_B=M(R)'/(M^{P(R)}(R)'\cap P M^{P(R)}(R)')|_R}$) is established in a 
similar way. In view of ${S_B\cong \nu(B) S}$, ${S=R, P(R), O(R), 
O(P(R)), I(P(R))}$, lemma 2.5 and the observations before Lemma 2.6, up to 
isomorphism ${P(R)_B=P(R_B), I(P(R))_B=I(P(R)_B)}$, 
\begin{multline*}
\End(\nu(B) P(R))_{M(R)'}\ =\ \CM(R)|_{\nu(B) P(R)}\ =
\\ 
\End(\nu(B) P(R))_{\nu(B) M(R)'}\cong 
\nu(B) \CM(R)\cong \CM(R_B)
\end{multline*}
and, as a consequence, ${O(R)_B=O(R_B), O(P(R))_B=O(P(R)_B)=O(P(R_B))}$,
\begin{multline*}
\End(\nu(B) O(P(R)))_{M(P(R))'}\ =\ 
\End(O(P(R)))_{M(P(R))'}|_{\nu(B) O(P(R))}\ =
\\ 
\CM(R) \Id_{\nu(B) O(P(R))}\ =\ 
\End(\nu(B) O(P(R)))_{\nu(B) M(P(R))'}\ =
\\  
\End(\nu(B) O(P(R)))_{M(O(P(R)))'}\ =\ 
\End(\nu(B) O(P(R)))_{\nu(B) M(O(P(R)))'}\cong
\\ 
\End(O(P(R_B)))_{M(P(R_B))'}\ =\ \End(O(P(R_B)))_{M(O(P(R_B)))'}\cong 
\nu(B) \CM(R)
\end{multline*}
(Proposition 2.4, p. 4; quasi-injectivity of $P(R)_{M(R)'}$ and 
$O(P(R))_{M(P(R))'}$, complete in-\linebreak variance of ${\nu(B) P(R)_{M(R)'}\subseteq 
P(R)}$ and ${\nu(B) O(P(R))_{M(P(R))'}\subseteq O(P(R))}$). 
\end{proof}

\end{document}
